%% P32 --- Transformer wyprowadzony od podstaw
%%
%% Bliźniak: programs/en/P32-transformer-derived.tex. Ramka za ramką, makro za
%% makrem: parity.py porównuje oba wydania W KOLEJNOŚCI, więc zdanie zbudowane
%% po polsku musi nieść te same spany matematyczne i te same literały liczbowe
%% w tej samej kolejności. Cyfra zostaje cyfrą, a słowo zostaje słowem.
%%
%% ODNIESIENIE PRZED MATEMATYKĄ. Angielskie \enquote{Program P25's divisor}
%% musi po polsku zaczynać się od odniesienia, bo C4 czyta kolejność.
%%
%% Uwaga tłumaczeniowa: \enquote{strumień} za \enquote{stream}, bo to zwykłe
%% polskie słowo. \enquote{transformer}, \enquote{softmax} i \enquote{token}
%% zostają po angielsku i odmieniają się po polsku --- tak mówią polscy
%% inżynierowie, a kalka byłaby słowem, którego nikt nie wyszuka.

\program{Transformer wyprowadzony od podstaw}
\label{prog:P32}
\index{transformer}

\begin{outcomes}
\outcome{1--6}{nazwać każdą operację w bloku i powiedzieć, który program tej
  książki już ci ją dał}
\outcome{7--12}{powiedzieć, dlaczego wynik dzieli się przez pierwiastek, i
  czy ten argument przeżywa to, że wyniki są wyliczane, a nie losowane}
\outcome{13--17}{powiedzieć, co softmax daje warstwie ponad zrobieniem liczb
  dodatnimi}
\outcome{18--20}{powiedzieć, dlaczego podział warstwy na głowice nic nie
  kosztuje}
\outcome{21--25}{powiedzieć, dlaczego połączenie rezydualne zmienia to, czym
  gradient jest, a nie tylko powiększa jego czynniki}
\outcome{26--28}{powiedzieć, gdzie pozycja wchodzi do obliczenia, które poza
  tym jest ślepe na kolejność}
\outcome{29--35}{policzyć parametry bloku, powiedzieć, która jego połowa ma
  ich więcej, i wyprowadzić cache z tych samych kształtów}
\outcome{36--39}{powiedzieć, które wyprowadzenia tej książki opierają się na
  hipotezie niesprawdzonej na wytrenowanym modelu}
\end{outcomes}

\begin{quiz}
\item Blok rzutuje strumień cztery razy, a potem uruchamia dwuwarstwową sieć
  gęstą. Która z tych dwóch połów ma więcej parametrów?
  \teachesat{29--30}
  \answerto{Sieć gęsta, i to $2$ razy: uwaga to $4d^{2}$, a warstwa gęsta to
  $8d^{2}$.}
\item Uwaga dzieli swoje wyniki przez $\sqrt{d_k}$. W bloku zapytania i
  klucze nie są losowane, tylko wyliczane ze strumienia. Czy argument za tym
  dzielnikiem nadal obowiązuje? \teachesat{8--9}
  \answerto{Tak, i to dokładnie: złożony wynik jest formą dwuliniową, której
  wariancja wychodzi $d_k$ w wartości oczekiwanej, więc nic się nie zmienia.}
\item Zwykły łańcuch $\val{p32.depth}$ warstw ma gradient ograniczony przez
  $\val{p32.plain.bound}$. Co robi z tym dodanie połączenia rezydualnego do
  każdej warstwy? \teachesat{23--24}
  \answerto{Sprawia, że gradient przestaje być pojedynczym iloczynem. Staje
  się sumą po ścieżkach, z których jedna jest identycznością i wnosi $1$.}
\item Model o $\val{p32.layers}$ warstwach szerokości $\val{p32.d}$ obsługuje
  rozmowę. Co rośnie szybciej wraz z długością sekwencji: arytmetyka czy
  cache? \teachesatone{35}
  \answerto{Arytmetyka, która jest kwadratowa; cache jest liniowy. To dwa
  wykładniki w jednej warstwie.}
\item Które z twierdzeń, jakie ta książka wyprowadziła o uwadze, sprawdzono
  na modelu naprawdę wytrenowanym? \teachesat{36--37}
  \answerto{Żadnego. Każde jest wyprowadzone przy inicjalizacji, a ten
  program mówi to wprost, zamiast sugerować coś innego.}
\end{quiz}

% ============================================================
\section{Jedna pozycja przez jeden blok}
\label{sec:P32-one-block}
\index{transformer!jeden blok}

\begin{fr}
To ostatnia część książki i nie wprowadza ona żadnej nowej matematyki.

Nie jest to skromność wobec materiału. To projekt: trzydzieści jeden
programów zbudowało po jednym kawałku, a zostało ułożyć je w kolejności i
zobaczyć, ile ten układ kosztuje.

Oto więc cały blok, jako sześć operacji wykonanych po kolei na wektorze
jednej pozycji: zrzutuj go na trzy sposoby; weź po jednym iloczynie skalarnym
na każdą parę pozycji; podziel je przez coś; zamień je na wagi sumujące się
do jedności; uśrednij trzecie rzutowanie tymi wagami; a potem dodaj wynik z
powrotem do tego, od czego zacząłeś, i przeskaluj.

Zanim pójdziesz dalej: ile z tych sześciu to operacje, które ta książka już
zdefiniowała?
\dotline
\nextframe
\end{fr}

\begin{fr}
\ans{Wszystkie sześć}

Każda, wraz z programem, który ci ją dał:

\begin{itemize}
\item \textbf{Zrzutuj na trzy sposoby.} Macierz jest funkcją
  (Program~\ref{prog:P06}), a trzy z nich zastosowane do jednego wektora to
  trzy funkcje zastosowane do jednego argumentu.
\item \textbf{Po jednym iloczynie skalarnym na parę.} Iloczyn skalarny
  (Program~\ref{prog:P05}), a Program~\ref{prog:F09} zdefiniował go
  wcześniej.
\item \textbf{Podziel przez coś.} Program~\ref{prog:P25} wyprowadza, przez co
  i dlaczego, a sekcja~\ref{sec:P32-score} jest o tym, czy to wyprowadzenie
  przeżywa złożenie.
\item \textbf{Zamień je na wagi sumujące się do jedności.} $\softmax$,
  którego stabilna postać należy do Programu~\ref{prog:P02}.
\item \textbf{Uśrednij tymi wagami.} Kombinacja wypukła
  (Program~\ref{prog:P19}).
\item \textbf{Dodaj z powrotem i przeskaluj.} Dwie połowy
  sekcji~\ref{sec:P32-residual} i sekcji~\ref{sec:P32-norm}.
\end{itemize}

W bloku nie ma nic więcej. Architekturę tworzy kolejność i kształty, a
kształty warto mieć dokładnie dobre, bo Program~\ref{prog:P07} poświęcił cały
program temu, co się dzieje, gdy nie są.

Strumień niesie na każdej pozycji wektor szerokości $\val{p32.d}$. Zapytania
powstają przez macierz $W_Q$. Jakiego kształtu musi być $W_Q$?
\dotline
\nextframe
\end{fr}

\begin{fr}
\ans{$\val{p32.d}$ na $d_k$}

Czyta strumień, więc bierze $\val{p32.d}$ liczb na wejściu; produkuje
zapytanie, więc wypuszcza $d_k$. Program~\ref{prog:P06} nazywa to funkcją z
jednej przestrzeni w drugą, a Program~\ref{prog:P07} nazywa to sygnaturą
typu; to jedno i to samo stwierdzenie.

$W_K$ ma ten sam kształt, bo klucz musi spotkać zapytanie w iloczynie
skalarnym, a Program~\ref{prog:F09} wymaga do tego dwóch wektorów tej samej
długości. $W_V$ prowadzi strumień do takiej szerokości, jaką mają wartości.

Zanim pójdziesz dalej: trzy z czterech rzutowań są już policzone. Po co w
ogóle istnieje czwarte? Wartości są przecież wektorami \dash{} czemu nie
dodać ich ważonej średniej wprost z powrotem?
\dotline
\nextframe
\end{fr}

\begin{fr}
\ans{Bo szerokości się nie zgadzają i bo głowice trzeba połączyć}

Wartości żyją na szerokości $d_k$, a strumień ma $\val{p32.d}$, więc coś musi
przenieść jedno w drugie; to sygnatura typu Programu~\ref{prog:P07}
odmawiająca dodawania, a nie wybór stylistyczny.

A gdy głowic jest kilka, ich wyjścia stoją obok siebie i trzeba je zmieszać.
Czwarte rzutowanie robi obie rzeczy naraz i dlatego jest jedną macierzą, a
nie dwiema.

\textbf{Cztery rzutowania i to całe parametry uwagi.} Zapamiętaj te cztery;
sekcja~\ref{sec:P32-count} je liczy.

Teraz część, która przesądza o wszystkim dalej. Wyniki powstają między
\emph{każdą} pozycją zapytania a \emph{każdą} pozycją klucza. Ile to wyników
na sekwencji $n$ pozycji?
\dotline
\nextframe
\end{fr}

\begin{fr}
\ans{$n^{2}$, i Program~\ref{prog:F10} je policzył}

Po jednym na parę uporządkowaną. Program~\ref{prog:P03} wycenił to i zrobił z
tego ostrzeżenie: parametrów jest $O(d^{2})$, a pracy $O(d\,n^{2})$, więc
liczba parametrów nie jest liczbą pracy.

To $n^{2}$ jest najbardziej brzemienną liczbą w tej architekturze. Każda
metoda z \enquote{długim kontekstem} w nazwie próbuje jej nie zapłacić, a
sekcja~\ref{sec:P32-count} stawia ją obok jedynej wielkości w bloku, która
\emph{nie} jest kwadratowa.

\begin{aibox}
Przeczytaj kartę modelu, a dostaniesz dwie liczby: liczbę parametrów i
długość kontekstu. Pierwsza wycenia macierze z tej sekcji, druga wycenia
wyniki. To różne wykładniki w tej samej warstwie, więc karta, która podwaja
kontekst, nie podwoiła niczego, co odczytasz z liczby parametrów \dash{}
poczwórnie powiększyła macierz wyników i podwoiła cache. Żadnej z tych liczb
na karcie nie ma.
\end{aibox}
\end{fr}

\mermaidfig{p32-one-block}{Jedna pozycja przez jeden blok}{sześć operacji}

\begin{fr}
Jeszcze jeden kształt, zanim zacznie się arytmetyka, bo to ten, który czyni
resztę programu czytelną.

Blok nie zastępuje strumienia. On do niego \emph{dodaje}: wektor, który
wchodzi, wychodzi z powrotem z dodanym wyjściem bloku. Strumień biegnie więc
od zanurzenia na dole do straty na górze, a każdy blok go czyta, coś liczy i
to coś dodaje z powrotem.

Dlatego nazywa się strumieniem, a nie potokiem, i
sekcja~\ref{sec:P32-residual} pokazuje, że wybór \emph{dodaj} zamiast
\emph{zastąp} jest powodem, dla którego głęboki stos w ogóle się trenuje.
\end{fr}

% ============================================================
\section{Wynik i to, czy wyprowadzenie przeżywa złożenie}
\label{sec:P32-score}
\index{uwaga!wynik}
\index{uwaga!dzielnik pierwiastkowy}

\begin{fr}
Wynik to jeden iloczyn skalarny: zapytanie z jednej pozycji przeciwko
kluczowi z innej. Program~\ref{prog:F09} go zdefiniował, a
Program~\ref{prog:P05} powiedział, co on mierzy.

I dzieli się go przez $\sqrt{d_k}$, zanim stanie się z nim cokolwiek innego.
Program~\ref{prog:P25} wyprowadził ten dzielnik, zamiast go stwierdzić, w
sekcji zbudowanej wokół wydobycia go od czytelnika, więc ten argument już raz
wyprodukowałeś.

Zanim pójdziesz dalej: powiedz, dlaczego dzielnik jest pierwiastkiem i
dlaczego akurat pierwiastkiem z szerokości głowicy.
\dotline
\nextframe
\end{fr}

\begin{fr}
\begin{ansblock}
Wynik to suma $d_k$ iloczynów. Wariancje wielkości niezależnych się dodają,
więc suma $d_k$ z nich ma wariancję $d_k$ i rozrzut $\sqrt{d_k}$. Dzielenie
przez ten pierwiastek zwraca wynik na skalę niezależną od $d_k$.
\end{ansblock}

Program~\ref{prog:P25} zmierzył, co się dzieje bez tego: przy
$d_k = \val{p25.e9.d.hi}$ rozrzut sięga $\val{p25.e9.raw.sd.512}$, jeden
klucz z $\val{p25.e9.keys}$ zabiera $\val{p25.e9.raw.top.512}$ procent wagi, a
własna pochodna $\softmax$ w tym punkcie wynosi $\val{p25.e9.resp.raw}$
przeciwko $\val{p25.e9.resp.scaled}$ z dzieleniem na miejscu.

Dzielnik jest więc poprawką wariancji, a nie stałą, którą ktoś dostroił. Nic
z tego nie jest do udowodnienia jeszcze raz przez ten program.

\textbf{Ale Program~\ref{prog:P25} zmierzył to na wektorach, które sam
wylosował}, powiedział to i przekazał tutaj jedną rzecz: w bloku zapytania i
klucze nie są losowane. Są \emph{wyliczane}, oba z tego samego strumienia,
przez dwie macierze, które wybrał trening.

To inne obliczenie i warto dokładnie powiedzieć, która jego część jest
podejrzana. Argument Programu~\ref{prog:P25} potrzebuje, żeby wyrazy
wchodzące do iloczynu skalarnego były od siebie niezależne. Tutaj są dwoma
liniowymi obrazami \emph{tego samego} strumienia, są więc funkcjami wspólnego
źródła, a \enquote{wylosowane niezależnie} to dokładnie to założenie, które
złożenie mogło zniszczyć.

Warto wiedzieć, co od tej odpowiedzi zależy, zanim się jej udzieli. Jeśli
założenie nie przetrwa, to dzielnik jest poprawny wobec modelu
Programu~\ref{prog:P25} i błędny wobec każdego bloku, jaki ktokolwiek
wytrenował, a ta sekcja musiałaby to powiedzieć, a nie powtarzać za nim.
Pytanie nie jest formalnością i jest jedynym w tym programie, którego
odpowiedzi książka nie ustaliła już gdzie indziej.

Zanim pójdziesz dalej: czy argument o wariancji nadal stosuje się do wyniku
policzonego w ten sposób?
\dotline
\nextframe
\end{fr}

\begin{fr}
\ans{Stosuje się i to dokładnie \dash{} ale nie z tego powodu, na jaki
wygląda}

Rozpisz wynik z rzutowaniami na miejscu. Zapytanie to $W_Q\T x$, a klucz to
$W_K\T y$, więc
\[ q \cdot k \;=\; (W_Q\T x) \cdot (W_K\T y) \;=\; x\T M y,
   \qquad M = W_Q W_K\T. \]

Warto się przy tym zatrzymać. \textbf{Dwie macierze rzutowania nigdy nie
występują w wyniku osobno.} Występuje tylko ich iloczyn, a ich iloczyn to
jedna macierz $\val{p32.d}$ na $\val{p32.d}$. Wynik jest formą dwuliniową:
jeden wektor, jedna macierz, drugi wektor.

A dla dwóch niezależnych strumieni o rozkładzie normalnym standardowym
wariancja $x\T M y$ wynosi dokładnie $\lVert M\rVert_F^{2}$ \dash{} sumę
kwadratów wszystkich wyrazów $M$. Do jej znalezienia nie trzeba żadnego
losowania.

To całe pytanie o złożenie, więc warto nazwać objazd, który ono zastępuje.

\begin{note}
Przed tym projektem próbowano dwóch innych i oba były złe, a ciekawsza jest
druga porażka.

Pierwszy losował świeże $W_Q$ i $W_K$ w każdej próbie. \textbf{To nie jest
transformer.} Blok ma ustalone wagi i zmienne wejścia, więc uśrednianie po
losowaniach wag odpowiada na łatwiejsze pytanie niż to zadane.

Drugi ustalił wagi i losował wejścia, i był po prostu za wolny. Wtedy
powyższa tożsamość uczyniła losowanie całkiem zbędnym \dash{} i to jest
odkrycie, a nie skrót, bo wielkość, którą umiesz zapisać dokładnie, nie
potrzebuje słupka błędu.
\end{note}
\end{fr}

\begin{fr}
Teraz wartość oczekiwana, i są to cztery linijki.

Reguła wachlarza wejściowego, którą wyprowadza Program~\ref{prog:P25}, kładzie
wariancję każdego wyrazu macierzy rzutowania na $1/\val{p32.d}$, tak by
zrzutowany wektor wychodził na tej skali, na jakiej wszedł. Wtedy każdy wyraz
$M$ jest sumą $d_k$ iloczynów dwóch takich wyrazów i

\[ \Ex\lVert M\rVert_F^{2}
   \;=\; \val{p32.d}^{2} \times d_k \times
         \Bigl(\frac{1}{\val{p32.d}}\Bigr)^{2}
   \;=\; d_k. \]

\textbf{Złożenie zachowuje wariancję dokładnie.} Dwa rzutowania między
strumieniem a wynikiem nic w odpowiedzi nie zmieniają, więc
Program~\ref{prog:P25} daje w $\sqrt{d_k}$ właściwy dzielnik dla prawdziwego
bloku, a nie tylko dla pary wylosowanych wektorów.

Zmierzone, przy $d_{\text{model}} = \val{p32.asm.dmodel}$ na
$\val{p32.asm.draws}$ losowaniach wag, rozrzut złożony przeciwko przewidywaniu:

\begin{center}
\begin{tabular}{rrr}
\toprule
$d_k$ & rozrzut złożony & $\sqrt{d_k}$ \\
\midrule
$8$  & $\val{p32.asm.sd.8}$  & $\num{2.83}$ \\
$16$ & $\val{p32.asm.sd.16}$ & $\num{4.00}$ \\
$32$ & $\val{p32.asm.sd.32}$ & $\num{5.66}$ \\
$64$ & $\val{p32.asm.sd.64}$ & $\num{8.00}$ \\
\bottomrule
\end{tabular}
\end{center}

Pierwszy i ostatni wiersz to własne szerokości głowic
Programu~\ref{prog:P25}, a jego zapisane tam rozrzuty to
$\val{p25.e9.raw.sd.8}$ i $\val{p25.e9.raw.sd.64}$. Skrypt sprawdza, że
złożone liczby tego programu je odtwarzają, więc oba programy nie mogą się
rozejść co do wielkości, którą wspólnie opisują.
\end{fr}

\begin{fr}
A teraz połowa, do której metoda Programu~\ref{prog:P25} nie mogła sięgnąć i
dla której przekazał on to pytanie dalej.

$\Ex\lVert M\rVert_F^{2} = d_k$ mówi o \emph{średniej} po losowaniach wag.
Wytrenowany model nie ma średniej. Ma jedno losowanie \dash{} jedno $M$
\dash{} a jego wyniki robią to, co $\lVert M\rVert_F^{2}$ dla tej jednej
macierzy.

Zmierzony po losowaniach względny rozrzut tej wielkości:

\begin{center}
\begin{tabular}{rr}
\toprule
$d_k$ & rozrzut $\lVert M\rVert_F^{2}$ \\
\midrule
$8$  & $\val{p32.asm.spread.8}$\,\% \\
$16$ & $\val{p32.asm.spread.16}$\,\% \\
$32$ & $\val{p32.asm.spread.32}$\,\% \\
$64$ & $\val{p32.asm.spread.64}$\,\% \\
\bottomrule
\end{tabular}
\end{center}

\textbf{To liczby orientacyjne, a nie dokładne}, i warto powiedzieć dlaczego:
$\val{p32.asm.draws}$ losowań szacuje rozrzut z dokładnością do jakiejś jego
siódmej części, więc ta kolumna jest rzędem wielkości, a nie pomiarem. Co nie
jest orientacyjne, to kolejność, a kolejność ma powód.

Warto jasno powiedzieć, co rozrzut po losowaniach znaczy dla czegoś, co się
wdraża. Dwa modele trenowane z różnych ziaren startują od różnych $M$, więc
ich wyniki siedzą na nieco innych skalach, zanim którykolwiek zobaczy
dane \dash{} i to jest jedno z miejsc, z których bierze się zmienność między
ziarnami w treningu, a nie cokolwiek, co zrobił optymalizator.

Zanim pójdziesz dalej: $M$ ma $\val{p32.d}^{2}$ wyrazów niezależnie od $d_k$.
Dlaczego \emph{wąska} głowica miałaby dawać bardziej zmienne
$\lVert M\rVert_F^{2}$ niż szeroka?
\dotline
\nextframe
\end{fr}

\begin{fr}
\ans{Bo rząd $M$ jest co najwyżej $d_k$}

$M = W_Q W_K\T$ jest iloczynem przez talię szerokości $d_k$, więc stosuje się
do niego wprost twierdzenie o wąskim gardle z Programu~\ref{prog:P08}: jego
rząd nie może przekroczyć $d_k$, jakkolwiek duża by ta macierz była.

Wyrazy $M$ nie są więc $\val{p32.d}^{2}$ niezależnymi liczbami. Są
$\val{p32.d}^{2}$ liczbami zbudowanymi z $d_k$ niezależnych kierunków, a wąska
głowica ma ich mniej do uśrednienia. Dlatego rozrzut maleje wraz z
poszerzaniem głowicy, i jest to Program~\ref{prog:P08} wykonujący zadanie,
którego nikt mu nie pisał.

\begin{trapbox}
\textbf{\enquote{Skalowanie jest dokładne, więc wyniki prawdziwego modelu
mają wariancję $d_k$.}}

Pierwszy człon jest prawdziwy, a drugi z niego nie wynika. Dokładność w
wartości oczekiwanej mówi o zbiorze modeli, które mogłeś wytrenować;
wdrożenie ma jeden, a jego wyniki są o kilka procent obok, zanim trening się
zacznie.

To rozróżnienie Programu~\ref{prog:P21} \dash{} nieobciążony to nie to samo
co użyteczny \dash{} przybywające do architektury zamiast do estymatora.
Efekt jest tu mały i warto wiedzieć, że nie jest zerowy, bo pytaniem następnej
sekcji jest, co się z nim dzieje po treningu, a na to nic w tej książce nie
odpowiada.
\end{trapbox}
\end{fr}

\mermaidfig{p32-score-assembled}{Skąd bierze się wynik}{jedna forma}

% ============================================================
\section{Od wyników do wag}
\label{sec:P32-softmax}
\index{transformer!softmax na wynikach}

\begin{fr}
Wyniki są teraz rzędem liczb rzeczywistych, po jednej na pozycję klucza, i
mogą mieć dowolny znak i dowolny rozmiar. Muszą stać się wagami.

Program~\ref{prog:F05} rozstrzygnął, co robi z nimi $\softmax$, a
Program~\ref{prog:P02} rozstrzygnął, jak go policzyć bez przepełnienia. Oba
są wydane; żadnego się tu nie powtarza.

Warto natomiast zapytać, co warstwa przez ten wybór \emph{zyskuje}, bo sama
dodatniość byłaby dużo tańsza. Podniesienie każdego wyniku do kwadratu daje
liczby dodatnie; samo eksponowanie bez niczego dalej też. Żadne z nich nie
sumuje się do jedności, a dzielenie na końcu $\softmax$-a jest jedynym
krokiem, który to zapewnia \dash{} pytanie nie brzmi więc, co daje
eksponenta, tylko co daje \emph{normalizacja}.

Wagi wychodzą dodatnie i sumują się do jedności. Co wobec tego można
powiedzieć o tym, gdzie musi leżeć wyjście warstwy względem wektorów
wartości, które uśrednia?
\dotline
\nextframe
\end{fr}

\begin{fr}
\ans{Wewnątrz ich powłoki wypukłej}

To dokładnie obiekt Programu~\ref{prog:P19}. Wagi nieujemne i sumujące się do
jedności tworzą kombinację wypukłą, więc wyjście jest ważoną średnią wartości
i nigdy nie może opuścić obszaru, który one rozpinają.

To prawdziwa gwarancja i warto ją nazwać, bo mówi, czego pojedyncza warstwa
uwagi strukturalnie \emph{nie potrafi}: nie potrafi wyprodukować kierunku, w
którym nie wskazuje żadna z jej wartości. Cokolwiek zapisze do strumienia,
jest zbudowane z tego, co już tam było.

Warstwa gęsta, która potem następuje, jest tym, co się z tego wymyka, bo
Program~\ref{prog:F05} udowodnił, że złożenie odwzorowań afinicznych jest
afiniczne, a nieliniowość jest tym, co to łamie. Dwie połowy bloku robią
różne rzeczy, i to jest zdanie, które mówi która.

\begin{aibox}
Własność powłoki wypukłej jest powodem, dla którego czysto uwagowy stos nie
może wymyślić cechy: na każdej warstwie wyjście siedzi wśród własnych wejść.
Gdy wynik interpretowalnościowy zgłasza kierunek pojawiający się na warstwie
$k$, a nieobecny na warstwie $k-1$, warstwa gęsta jest tym, skąd on mógł
przyjść, a Program~\ref{prog:P31} jest programem o tym, na co taki wynik
pozwala, a na co nie.
\end{aibox}
\end{fr}

\begin{fr}
Istnieje tańszy sposób, żeby liczby sumowały się do jedności, i warto
zobaczyć, dlaczego się go nie używa: podzielić każdy wynik przez sumę
wyników.

Zanim pójdziesz dalej: wyniki mogą w tym miejscu mieć dowolny znak. Co to
robi tej tańszej metodzie?
\dotline
\nextframe
\end{fr}

\begin{fr}
\ans{Psuje ją, i to na dwa osobne sposoby}

Ujemny wynik dałby ujemną wagę, a wyjście przestałoby w ogóle być średnią
wartości \dash{} gwarancja powłoki wypukłej znika. Gorzej, wyniki mogą sumować
się do zera i wtedy metoda dzieli przez zero z powodów niemających nic
wspólnego z tym, że dane są nietypowe.

Wcześniejsze potęgowanie naprawia oba: posyła każdy wynik w liczbę dodatnią, a
suma liczb dodatnich nigdy nie jest zerem. \textbf{Eksponenta jest po to, by
dzielenie było bezpieczne, a gwarancja prawdziwa}, a to, że przy okazji
zamienia różnice wyników na ilorazy wag, pokazał Program~\ref{prog:F07} i to
właśnie czyni dzielnik z sekcji~\ref{sec:P32-score} tak ważnym.
\end{fr}

\begin{fr}
Jedna przestroga co do wag, i należy do Programu~\ref{prog:P29}.

Wagę $\num{0.9}$ na jednym kluczu czyta się często jako
\enquote{skupienie uwagi} warstwy na tym kluczu, a ten odczyt mówi o
rozkładzie, więc da się go zmierzyć zamiast oceniać na oko:
Program~\ref{prog:P29} zamienia entropię wiersza uwagi na efektywną liczbę
kluczy i podaje $\val{p29.att.raw.eff}$ z $\val{p29.att.keys}$ bez dzielnika
przeciwko $\val{p29.att.scaled.eff}$ z nim.

Dzielnik z sekcji~\ref{sec:P32-score} nie jest więc numerycznym drobiazgiem.
Jest tym, co powstrzymuje warstwę przed zapadnięciem się na jeden klucz,
zanim się czegokolwiek nauczy, a jego efekt daje się odczytać w jednostce, w
której spiera się przegląd projektowy.
\end{fr}

% ============================================================
\section{Głowice}
\label{sec:P32-heads}
\index{transformer!wiele głowic}

\begin{fr}
Blok nie uruchamia jednej uwagi, tylko kilka równolegle i skleja wyniki.
Program~\ref{prog:P07} ma mechanikę \dash{} szeroki wektor dzieli się na $h$
plastrów, oś głowicy przenosi się przed oś pozycji, a jedno zwężenie robi
resztę \dash{} i nazywa to przekształceniem kształtu, a nie obliczeniem.

Oto ile to kosztuje. Przy $\val{p32.heads}$ głowicach szerokości
$\val{p32.d.head}$ szerokości głowic sumują się do
$\val{p32.heads} \times \val{p32.d.head} = \val{p32.d}$, czyli do własnej
szerokości strumienia.

Zanim pójdziesz dalej: porównaj parametry w $\val{p32.heads}$ głowicach
szerokości $\val{p32.d.head}$ z jedną głowicą szerokości $\val{p32.d}$.
\dotline
\nextframe
\end{fr}

\begin{fr}
\ans{Są takie same, dokładnie}

$W_Q$ dla jednej szerokiej głowicy ma wymiar $\val{p32.d}$ na
$\val{p32.d}$. Dla $\val{p32.heads}$ wąskich głowic jest to
$\val{p32.heads}$ macierzy $\val{p32.d}$ na $\val{p32.d.head}$ \dash{} a
$\val{p32.heads} \times \val{p32.d.head} = \val{p32.d}$, więc kolumny składają
się na tę samą macierz.

\textbf{Uwaga wielogłowicowa jest więc darmowa.} To jedno rzutowanie
$\val{p32.d}$ na $\val{p32.d}$ czytane jako $\val{p32.heads}$ bloków zamiast
jednego, i dlatego Program~\ref{prog:P07} mógł nazwać to przekształceniem
kształtu: nic nie liczy się inaczej, liczby są inaczej pogrupowane.

Zmienia się nie koszt, lecz ograniczenie. Jedna szeroka głowica tworzy jedno
$M = W_Q W_K\T$ rzędu do $\val{p32.d}$; $\val{p32.heads}$ wąskich głowic
tworzy $\val{p32.heads}$ osobnych, każde rzędu co najwyżej
$\val{p32.d.head}$. Blok wymienia jedno bogate porównanie na wiele
ograniczonych przy identycznym koszcie \dash{} a Program~\ref{prog:P08} jest
miejscem, gdzie słowo \enquote{ograniczone} zostało uściślone.

Program~\ref{prog:P07} stawia jedną tezę o tym układzie, którą łatwo pominąć.
Oś głowicy przenosi się \emph{przed} oś pozycji, zanim nastąpi zwężenie.

Zanim pójdziesz dalej: co poszłoby źle, gdyby tak nie było?
\dotline
\nextframe
\end{fr}

\begin{fr}
\ans{Spotkałyby się pozycje z różnych głowic}

Zwężenie biegnie po ostatniej osi i paruje niezależnie wszystko, co stoi
przed nią. Z osią głowicy z przodu każda głowica porównuje własne zapytania z
własnymi kluczami i głowice nigdy się nie mieszają. Przy odwrotnej kolejności
osi to samo zwężenie sparowałoby zapytanie z jednej głowicy z kluczem z
innej, co nie jest porównaniem czegokolwiek.

Żaden błąd by nie powstał. Kształty wychodzą tak czy inaczej, a
Program~\ref{prog:P07} jest całym programem o tym, że to jest przypadek
niebezpieczny \dash{} układ, który działa i liczy coś innego, niż zamierzono.

\begin{trapbox}
\textbf{\enquote{Więcej głowic to większa pojemność.}}

Więcej głowic przy ustalonej szerokości \dash{} owszem. Więcej głowic przy
ustalonej szerokości \emph{łącznej} \dash{} czyli to, co robi każda
architektura \dash{} to te same parametry pokrojone na węższe kawałki
niższego rzędu. Czy ta wymiana się opłaca, jest pytaniem empirycznym o klasę
zadań i ta książka na nie nie odpowiada; arytmetyka rozstrzyga tyle, że jest
to wymiana, a nie zysk.
\end{trapbox}
\end{fr}

% ============================================================
\section{Strumień rezydualny}
\label{sec:P32-residual}
\index{transformer!strumień rezydualny}
\index{gradient!a połączenia rezydualne}

\begin{fr}
Program~\ref{prog:F12} udowodnił coś o głębokości, a potem odmówił dokończenia
tego.

Jego argument: gradient przez stos warstw jest \emph{iloczynem}, jeden
czynnik na warstwę, a jeśli każdy czynnik jest co najwyżej $\frac{1}{4}$, to
całość jest co najwyżej $\bigl(\frac{1}{4}\bigr)^{n}$.

Zanim pójdziesz dalej: przy $\val{p32.depth}$ warstwach ile wynosi to
ograniczenie?
\dotline
\nextframe
\end{fr}

\begin{fr}
\ans{$\val{p32.plain.bound}$, i to w najlepszym przypadku}

Program~\ref{prog:F12} zapisuje dokładnie tę liczbę, a ten program przelicza
ją na tym samym łańcuchu, więc jeśli któreś się ruszy, budowanie staje.

Warto jasno powiedzieć, jak źle to jest, a także jak niezupełnie źle w ten
sposób, w jaki się mówi. Przy $\val{p32.plain.bound}$ liczba jest wciąż w
pełni reprezentowalna: Program~\ref{prog:P01} kładzie najmniejszą liczbę
normalną $\code{float32}$ na $\val{p01.fp32.minnorm}$, więc w najlepszym
przypadku nic się nie zapada.

Zapada się przypadek realistyczny. Program~\ref{prog:F12} liczy ten sam
iloczyn na warstwie \emph{nasyconej}, a nie w jej najlepszym punkcie, i
dostaje iloczyn poniżej $10^{-\val{f12.sat.exponent}}$ \dash{} pod każdą
podłogą, jaką ten format ma. Wtedy wczesne warstwy dostają dokładnie zero i
dostają je po cichu.

Program~\ref{prog:F12} nazywa następnie odpowiedź architektury i przekazuje
wyprowadzenie tutaj: \emph{ścieżka od straty do wczesnych warstw, która wcale
nie jest długim iloczynem.}

Oto więc zmiana, i jest to jeden symbol. Zamiast warstwy zastępującej
strumień, $y = f(x)$, warstwa do niego dodaje:
\[ y \;=\; x + f(x). \]

Zanim to zróżniczkujesz, warto zauważyć dwie rzeczy. To nakłada na blok
warunek: $x$ i $f(x)$ muszą mieć ten sam kształt, żeby w ogóle dało się je
dodać, i dlatego każda szerokość w
sekcji~\ref{sec:P32-one-block} wracała do $\val{p32.d}$. I nie wprowadza to
żadnej operacji, której ta książka by nie miała \dash{} to dodawanie, a
Program~\ref{prog:F11} podaje pochodną sumy.

Nic tu więc nie wymaga nowej reguły, a zmiana, która nie wymaga nowej reguły,
jest dokładnie tą, którą czyta się jako kosmetyczną. Warto samemu rozstrzygnąć,
co ona robi, zanim usłyszysz, po co jest.

Zanim pójdziesz dalej: zróżniczkuj to. Ile wynosi jakobian jednej takiej
warstwy?
\dotline
\nextframe
\end{fr}

\begin{fr}
\ans{$1 + f'(x)$}

Co wygląda na małą zmianę i nią nie jest. Ułóż $n$ takich warstw, a gradient
wynosi
\[ \prod_{k=1}^{n} \bigl(1 + f'_k\bigr), \]
podczas gdy wcześniej był $\prod_k f'_k$. Narzucający się odczyt jest taki, że
czynniki zrobiły się większe, i ten odczyt jest dostępny oraz błędny.

Wymnóż nawiasy. Każdy składnik bierze z każdego nawiasu albo $1$, albo
$f'_k$ \dash{} więc rozwinięcie ma jeden składnik na każdy \emph{podzbiór}
warstw, czyli $2^{n}$ składników, a każdy jest iloczynem $f'$ po warstwach ze
swojego podzbioru.

\textbf{To jest suma po ścieżkach.} Podzbiór mówi, przez które warstwy
gradient \emph{przeszedł}, a które \emph{ominął}.

Zauważ, co samo rozwinięcie zrobiło z pytaniem. Iloczyn to jedna wielkość i
jeden mały czynnik gdziekolwiek w nim rozstrzyga całość. Suma to $2^{n}$
wielkości i rozstrzyga ją ta największa, więc mały składnik nic nie kosztuje.
Oba wyrażenia są równe i nie jest jednakowo łatwo o nich myśleć \dash{} i
dlatego rozpisanie było warte tej linijki.

Zanim pójdziesz dalej: jeden z tych $2^{n}$ podzbiorów jest pusty. Ile wnosi
jego składnik?
\dotline
\nextframe
\end{fr}

\begin{fr}
\ans{Dokładnie $1$}

Pusty podzbiór bierze $1$ z każdego nawiasu. To ścieżka, która omija każdą
warstwę \dash{} prosto od straty do wejścia, nie dotykając niczego.

I to jest cała rzecz:

\begin{center}
\textbf{gradient zwykłego łańcucha to jeden iloczyn, który zabija $n$ małych
czynników.\\
Gradient stosu rezydualnego to suma, a jednym z jej składników jest $1$.}
\end{center}

Żaden czynnik go nie zabije, bo nie zawiera on żadnych czynników. Pozostałe
$2^{n} - 1$ ścieżek może znikać dokładnie tak, jak mówi
Program~\ref{prog:F12}, a gradient i tak dojdzie.

Skrypt sprawdza to rozwinięcie na ułamkach przy $\val{p32.path.n}$ warstwach
\dash{} $\val{p32.path.count}$ ścieżek sumujących się dokładnie do iloczynu
\dash{} a potem sprawdza konsekwencję, która jest tezą, a nie arytmetyką:
wyłącz każdą warstwę, tak by każde $f'_k$ było zerem, a gradient zwykłego
łańcucha wynosi dokładnie $0$, gradient stosu rezydualnego zaś dokładnie $1$.
\end{fr}

\transcript{p32-two-stacks}

\begin{fr}
Druga linia tego listingu to przypadek uczciwy, w którym warstwy robią coś
małego, a nie nic.

Na $\val{p32.res.trials}$ losowaniach, przy każdym $f'$ rozmiaru około
$\val{p32.res.sd}$ i $\val{p32.depth}$ warstwach, mediana gradientu zwykłego
łańcucha wynosi $\val{p32.res.plain.med}$, a stosu rezydualnego
$\val{p32.res.resid.med}$, przy czym $\val{p32.res.within}$ procent losowań
mieści się w granicach dwukrotności $1$.

Te dwie liczby są ilustracją dokładnego wyniku, a nie wynikiem, więc skrypt
sprawdza kolejność i to, że stos rezydualny zostaje rzędu jedności, a nigdy
samych liczb.

\begin{trapbox}
\textbf{\enquote{Połączenia rezydualne naprawiają znikające gradienty.}}

One zmieniają to, czym gradient \emph{jest}. To mocniejsze niż naprawienie i
zarazem węższe, a obie połowy mają znaczenie.

Mocniejsze: ścieżka identycznościowa nie jest dużym czynnikiem, który wciąż
może zostać przeważony, tylko składnikiem, którego nic nie mnoży.

Węższe: nic tutaj nie mówi, że pozostałe $2^{n} - 1$ ścieżek się zachowuje.
Stos rezydualny wciąż może eksplodować \dash{} druga awaria z
Programu~\ref{prog:F12} pozostaje nietknięta, i po to jest następna sekcja
\dash{} a dodanie połączeń rezydualnych do źle wyskalowanej sieci daje
gradient równy $1$ plus bardzo dużo szumu. Ścieżka identycznościowa gwarantuje
dojście, a nie przydatność.
\end{trapbox}
\end{fr}

\mermaidfig{p32-paths-not-product}{Dwa stosy, dwa gradienty}{suma i iloczyn}

% ============================================================
\section{Normalizacja i miejsce, w którym wchodzi pozycja}
\label{sec:P32-norm}
\index{transformer!normalizacja warstwy}
\index{transformer!informacja o pozycji}

\begin{fr}
Zostały dwie operacje i obie są wydane, więc ta sekcja mówi, gdzie one siedzą,
a nie czym są.

\textbf{Normalizacja} odejmuje średnią wektora i dzieli przez jego rozrzut.
Program~\ref{prog:P18} wyprowadza jej gradient i sprawdza go z różnicą
centralną z dokładnością lepszą niż $\val{p18.check.ceiling}$, więc nic o niej
nie jest tu należne.

Złożenie dodaje odpowiedź na to, po co ona w bloku w ogóle jest.
Sekcja~\ref{sec:P32-residual} zostawiła otwartą drugą awarię: ścieżka
identycznościowa gwarantuje, że gradient dojdzie, i nie mówi nic o rozmiarze
tego, co dojdzie razem z nim. Każdy blok dodaje do strumienia, więc skala
strumienia rośnie z głębokością, o ile coś jej nie trzyma.

Zanim pójdziesz dalej: co normalizacja robi z argumentem z
sekcji~\ref{sec:P32-score}?
\dotline
\nextframe
\end{fr}

\begin{fr}
\ans{Przywraca warunek, który ten argument zakłada}

Rachunek wariancji zakładał, że wyrazy strumienia mają znaną, ustaloną skalę.
Trzydzieści dwa bloki dodające do niego nie zostawiłyby tego prawdą.
Normalizacja odkłada to na miejsce, zanim każdy blok zacznie czytać, więc
wariancja wyniku trzyma się na warstwie $\val{p32.layers}$ z tego samego
powodu, z którego trzyma się na pierwszej.

Trzy kawałki są więc jednym mechanizmem, a nie trzema: ścieżka rezydualna
niesie gradient, normalizacja trzyma ustaloną skalę, na której jest on
niesiony, a dzielnik trzyma wyniki na tej skali.

\textbf{Informacja o pozycji} jest drugą i naprawia coś, czego blok w
dotychczasowej postaci po prostu nie ma.

Zanim pójdziesz dalej: weź blok z sekcji~\ref{sec:P32-one-block} i pomieszaj
pozycje wejściowe. Co się dzieje z jego wyjściami?
\dotline
\nextframe
\end{fr}

\begin{fr}
\ans{Mieszają się razem z nimi i poza tym są identyczne}

Każdy wynik jest iloczynem skalarnym dwóch wektorów, a ważona suma biegnie po
pozycjach, nie pytając o nie, więc nic w bloku nie odróżni \enquote{kot
usiadł tu} od \enquote{usiadł tu kot}. Nie jest tak, że kolejność jest
obsłużona źle; kolejność nie jest w ogóle reprezentowana.

Program~\ref{prog:F08} dostarcza naprawy i dowodzi własności, na której ona
stoi: obrócenie zapytania i klucza o kąty $\alpha$ i $\beta$ zostawia ich
iloczyn skalarny zależny tylko od $\beta - \alpha$. Obrócenie wektora każdej
pozycji o kąt proporcjonalny do jej indeksu czyni więc każdy wynik funkcją
\emph{odległości} między dwiema pozycjami i niczego więcej o tym, gdzie one
siedzą.

Obrót jest ortogonalny, więc nie zmienia żadnej długości \dash{}
Program~\ref{prog:F08} mierzy błąd na $\val{f08.rot.len.err}$, a
Program~\ref{prog:P09} mówi, dlaczego ortogonalna zmiana bazy nic model nie
kosztuje. \textbf{Pozycja zostaje dodana bez wydawania pojemności
strumienia}, i to jest cały powód, dla którego tej metody się używa.
\end{fr}

% ============================================================
\section{Ile to kosztuje}
\label{sec:P32-count}
\index{transformer!liczba parametrów}
\index{koszt!jednego bloku}

\begin{fr}
Program~\ref{prog:F01} otworzył tę książkę, każąc ci policzyć, ile waży model
o siedmiu miliardach parametrów. Nigdy nie powiedział, skąd bierze się te
siedem miliardów, a na trzeciej stronie nie było na to słownika.

Teraz jest. Program~\ref{prog:P03} zapisuje kształt tego samego modelu
\dash{} $\val{p32.layers}$ warstw szerokości $\val{p32.d}$ \dash{} a
sekcja~\ref{sec:P32-one-block} policzyła macierze.

Uwaga ma cztery rzutowania, każde $\val{p32.d}$ na $\val{p32.d}$, gdy głowice
są już połączone, więc mieści $4d^{2}$.

Warstwa gęsta to dwie warstwy, $d$ na wejściu i $4d$ na wyjściu, potem $4d$ na
wejściu i $d$ na wyjściu. Program~\ref{prog:F02} zebrał już te wyrazy
podobne: $4d^{2} + 4d^{2} = 8d^{2}$.

Oba rachunki są na jeden blok, więc żaden z nich nie zależy od głębokości:
pomnożenie przez liczbę warstw mnoży oba przez to samo i zostawia ich stosunek
dokładnie tam, gdzie był. I żaden nie jest nowy \dash{} pierwszy to cztery
rzutowania jednej szerokości, a drugi to zebrane wyrazy podobne
Programu~\ref{prog:F02}, porównanie jest więc między dwiema liczbami, które
ta książka już podała.

Oba liczą też macierze i nic poza nimi. Blok niesie jeszcze obciążenia i dwa
zestawy parametrów normalizacji, a tych jest po $\val{p32.d}$ liczb wobec
macierzy o $\val{p32.d}^{2}$ \dash{} mniej o czynnik równy szerokości, i
dlatego nikt ich nie liczy, i dlatego warto raz powiedzieć, że zostały
pominięte, a nie przeoczone.

Zanim pójdziesz dalej: która z dwóch połów bloku mieści więcej parametrów i o
ile?
\dotline
\nextframe
\end{fr}

\begin{fr}
\ans{Warstwa gęsta, i to $2$ razy}

$8d^{2}$ przeciwko $4d^{2}$. \textbf{Połowa, o której nikt nie mówi, mieści
dwa razy tyle, co połowa, o której mówią wszyscy.}

Warto przy tym posiedzieć, bo prawie każde zdanie napisane o tej
architekturze dotyczy uwagi, a dwie trzecie parametrów bloku siedzi w części,
która uwagą wcale nie jest. W liczbach, przy $d = \val{p32.d}$:

\begin{center}
\begin{tabular}{lr}
\toprule
uwaga, $4d^{2}$ & $\val{p32.attn.block}$ \\
warstwa gęsta, $8d^{2}$ & $\val{p32.ffn.block}$ \\
\midrule
jeden blok, $12d^{2}$ & $\val{p32.per.block}$ \\
$\times\ \val{p32.layers}$ warstw & $\val{p32.blocks.total}$ \\
\bottomrule
\end{tabular}
\end{center}

A teraz arytmetyka, której Program~\ref{prog:F01} nie mógł wykonać.

Zauważ, czym są te dwie liczby. Ta niżej to figura z karty modelu, czyli to,
co każdy by sprawdził; ta wyżej to, co wymuszają kształty, i nikt jej nie
wybierał. Cokolwiek dzieli te dwie liczby, musi się więc rozliczyć czymś, co
ten rachunek pominął, i warto zdecydować, czy spodziewasz się zaokrąglenia,
czy całego składnika.

Zanim pójdziesz dalej: $\val{p32.blocks.total}$ przeciwko liczbie, którą
podaje Program~\ref{prog:F01}, czyli $\val{p32.f01.params}$ \dash{} jaką
część tego modelu stanowią bloki?
\dotline
\nextframe
\end{fr}

\begin{fr}
\ans{$\val{p32.blocks.pct}$ procent}

\textbf{Bloki są więc modelem.} Reszta to zanurzenie, które przy słowniku
$\val{p32.vocab}$ daje $\val{p32.embed}$ parametrów, czyli
$\val{p32.embed.pct}$ procent \dash{} a dalsza reszta to głowica wyjściowa,
obciążenia i skale normalizacji, liczone w tysiącach tam, gdzie tamte liczy
się w miliardach.

Pierwszy program książki kazał ci zważyć liczbę, której nie umiał wyprowadzić.
Trzydzieści programów później liczba ta wychodzi z dwóch zapisanych kształtów
i jednego mnożenia.
\end{fr}

\begin{fr}
Dwa uczciwe zastrzeżenia do tego rachunku, bo to liczba tego rodzaju, który
podróżuje.

\begin{warning}
\textbf{To $12d^{2}$ należy do bloku oryginalnego, a nie do każdego.} Połowa
gęsta wynosi $8d^{2}$ tylko wtedy, gdy jej szerokość wewnętrzna to $4d$ i
używa ona dwóch macierzy. Kilka obecnych architektur używa trzech macierzy i
szerokości wewnętrznej innej niż $4d$, co zmienia $8$, a więc i $12$. Metoda
przeżywa \dash{} policz macierze, dodaj kwadraty \dash{} a stała nie.

A słownik jest \emph{wyborem projektowym}, a nie pomiarem.
Program~\ref{prog:F01} nigdy żadnego nie ustalił, więc $\val{p32.vocab}$ jest
tu nazwany jako założenie, którym jest; model o czterokrotnie większym
słowniku rusza wiersz zanurzenia, a nie wiersze bloków.
\end{warning}
\end{fr}

\begin{fr}
Teraz cache, a wychodzi on z tych samych kształtów.

Obsługa rozmowy oznacza nieprzeliczanie kluczy i wartości każdej wcześniejszej
pozycji na każdym kroku, więc trzyma się je. Trzymać trzeba jeden klucz i
jedną wartość na pozycję, na warstwę.

Zauważ, które z czterech rzutowań to są, a które nie. Zapytań się nie trzyma:
na każdym kroku jest dokładnie jedno nowe zapytanie, używa się go raz
przeciwko każdemu kluczowi i nigdy więcej nie jest potrzebne. Klucze i
wartości są tym, czego będzie potrzebował każdy przyszły krok. \textbf{Dwa z
czterech się przechowuje, a dwóch nie}, a które dwa \dash{} wynika z tego, co
obliczenie z nimi robi, a nie z czyjegoś wyboru.

Zanim pójdziesz dalej: przy $\val{p32.layers}$ warstwach szerokości
$\val{p32.d}$, po dwa bajty na liczbę, ile to jest na token?
\dotline
\nextframe
\end{fr}

\begin{fr}
\ans{$2 \times \val{p32.layers} \times \val{p32.d} \times \val{p32.bytes}$
bajtów, czyli $\val{p32.kv.per.token.mib}$ MiB}

Dwa tensory, po jednym wektorze $d$ każdy, na warstwę, na pozycję. Na
$\val{p32.seq}$ tokenach daje to $\val{p32.kv.gib}$ GiB \dash{} czyli własną
zapisaną liczbę Programu~\ref{prog:P03}, a skrypt sprawdza, że wyprowadzenie
tego programu z kształtów bloku ją odtwarza.

To złożenie zarabiające na siebie. Program~\ref{prog:P03} zmierzył cache; nie
mógł powiedzieć, dlaczego ma on taki rozmiar, bo powiedzenie tego wymaga
czterech rzutowań i faktu, że przechowuje się dokładnie dwa z nich. Teraz
może.

\begin{note}
Program~\ref{prog:P03} drukował tę liczbę na token jako \enquote{MB}, licząc
ją w mebibajtach, i przetrwało to cały szkic, bo
$\val{p32.kv.per.token.mib}$ MiB oraz $\num{0.524}$ MB zaokrąglają się do
$\val{p32.kv.per.token.mib}$ na jednym miejscu po przecinku. Jednostka jest
tam teraz poprawiona. To dokładnie ta pomyłka, przed którą ostrzega własne
podsumowanie Programu~\ref{prog:P03} \dash{} plan pojemności rozjechany o
siedem procent to prawie zawsze ta jedna \dash{} a powód, dla którego się
ukryła, należy do Programu~\ref{prog:P17}: dwa odczyty wzoru zgodne liczbowo
zostają niewidoczne, dopóki nie przestaną być.
\end{note}
\end{fr}

\begin{fr}
Ostatnia arytmetyka w programie jest tą, która rozstrzyga, jak architekturę się
wdraża.

Na $\val{p32.seq}$ pozycjach przy $\val{p32.heads}$ głowicach wyniki to
$\val{p32.scores}$ liczb. Cache to $\val{p32.cache.vectors}$ wektorów.

Podwój sekwencję, a pierwsze rośnie \emph{czterokrotnie}, drugie zaś
\emph{dwukrotnie} \dash{} dwa wykładniki w jednej warstwie, co jest odkryciem
Programu~\ref{prog:P03}, teraz odczytanym z kształtów, a nie stwierdzonym o
uwadze w ogólności.

Ten kwadratowy jest tym, którego trzeba się trzymać. \textbf{Każda metoda z
\enquote{długim kontekstem} w nazwie jest próbą niepłacenia $n^{2}$}:
liczenia tylko części par, liczenia ich z niższą precyzją, liczenia ich
blokami mieszczącymi się w szybkiej pamięci albo zastąpienia porównania
czymś, co nie jest porównaniem każdej pary. Różnią się tym, z czego rezygnują,
a rezygnują z tego, które pary wolno sobie zobaczyć.

\begin{aibox}
Dwa wykładniki tłumaczą wzorzec, który spotkał każdy, kto tę architekturę
obsługiwał. Model odpowiadający wygodnie na krótki prompt i wywracający się na
długim zwykle nie wyczerpał \emph{parametrów}: wyczerpał jedno z tych dwóch.
Które \dash{} mówi, co robić, a wskazują one w przeciwne strony: cache naprawia
się obsługą mniejszej liczby rozmów naraz, macierz wyników \dash{} obsługą
krótszych. Plan pojemności, który ich nie rozdziela, zgaduje.
\end{aibox}
\end{fr}

\mermaidfig{p32-where-parameters-are}{Gdzie siedzą parametry bloku}{cztery i osiem}

% ============================================================
\section{Czego ta książka nie sprawdziła}
\label{sec:P32-unchecked}
\index{pomiar!co pozostaje}

\begin{fr}
Architektura jest złożona, a każda liczba w niej wzięła się z programu, który
ją wyprowadził. Ta ostatnia sekcja jest o tych, które wyprowadzono tam, gdzie
model nie mieszka.

Sekcja~\ref{sec:P32-score} udowodniła, że złożenie zachowuje wariancję wyniku.
Przeczytaj ten dowód jeszcze raz i zapytaj, co on zakładał o strumieniu.
\dotline
\nextframe
\end{fr}

\begin{fr}
\ans{Że jego wyrazy były niezależne i wszystkie jednej skali}

Co jest prawdą przy inicjalizacji, bo inicjalizator czyni to prawdą.
Program~\ref{prog:P25} mówi to wprost i dodaje część, która ma znaczenie:
staje się to coraz mniej prawdą w miarę treningu, bo zapytania i klucze uczą
się struktury, a struktura to zależność.

Każde więc wyprowadzenie tej książki o uwadze \dash{} $\sqrt{d_k}$, wariancja
przez rzutowania, reguła wachlarza wejściowego, skala przywracana przez
normalizację \dash{} jest stwierdzeniem o modelu \textbf{przed jego
wytrenowaniem}.

\begin{warning}
\textbf{Ta książka nie zmierzyła żadnego z nich na modelu wytrenowanym i nie
ma takiego modelu.}

Uczciwe stanowisko jest takie, że wyprowadzenia są dokładne, a ich hipoteza
jest sprawdzalna tylko tam, dokąd ta książka nie sięga. Rozstrzygnęłoby to coś
nietrudnego i niemożliwego dla piaskownicy: wziąć wytrenowany model, policzyć
$\lVert W_Q W_K\T\rVert_F^{2}$ na głowicę i porównać z $d_k$. Jeśli te dwie
rzeczy się rozeszły, dzielnik nie robi już tego, do czego go wyprowadzono, a
każde dalsze twierdzenie o entropii wiersza uwagi dotyczy rozkładu, którego
skali nikt nie sprawdził.

Matematyka przewiduje, że się rozchodzą \emph{i że to nie ma znaczenia}, bo
$\sqrt{d_k}$ zostaje właściwym rzędem wielkości niezależnie od tego, co robi
zależność. To przewidywanie jest osądem. Jest tu oznaczone jako osąd, a nie
podane jako wynik.
\end{warning}
\end{fr}

\begin{fr}
Nie jest jedyne, a listę warto mieć w jednym miejscu, bo każda pozycja nazywa
program, który sam podał swoją granicę, zamiast ją zamalować.

\begin{center}
\begin{tabularx}{\linewidth}{lX}
\toprule
Program & Czego nie zmierzono \\
\midrule
\ref{prog:P08}, \ref{prog:P11} & czy widmo prawdziwej macierzy zanurzeń opada
  tak, jak zakłada argument o niskim rzędzie \\
\ref{prog:P19} & czy to, do której doliny trafi spacer, ma znaczenie w skali,
  w jakiej ludzie trenują \\
\ref{prog:P20} & czy metoda adaptacyjna dochodzi do lepszej odpowiedzi, a nie
  tylko dochodzi szybciej \\
\ref{prog:P25}, \ref{prog:P32} & czy niezależność zakładana przez
  wyprowadzenie wyniku przeżywa trening \\
\ref{prog:P31} & na co pozwala wynik sondy wobec warstwy, gdy zmierzy się
  lukę między ograniczeniem a prawdą \\
\bottomrule
\end{tabularx}
\end{center}

\textbf{Pięć pozycji i jedna przyczyna.} Każda potrzebuje prawdziwych macierzy
wytrenowanego modelu, a ta książka powstała tam, gdzie takiego nie ma. To
ograniczenie książki, a nie matematyki, a powód, dla którego jest ono tutaj
wydrukowane, a nie zostawione domyślnie, to dyscyplina, na której stoją oba
tomy towarzyszące: twierdzenie potrzebuje metody i liczby, albo zostaje
oznaczone jako osąd.
\end{fr}

\begin{fr}
Jedna uwaga na zamknięcie, i dotyczy tego, co pokazało złożenie, a nie samej
architektury.

Nic w tym programie nie było nowe. Wynik to iloczyn skalarny
Programu~\ref{prog:P05}; jego dzielnik to argument o wariancji
Programu~\ref{prog:P25}; wagi to kombinacja wypukła Programu~\ref{prog:P19};
głowice to przekształcenie kształtu Programu~\ref{prog:P07}; ścieżka
rezydualna to iloczyn Programu~\ref{prog:F12} przeczytany jako suma; rachunek
to zbieranie wyrazów podobnych Programu~\ref{prog:F02} wykonane dwanaście
razy; cache to pomiar Programu~\ref{prog:P03} wyprowadzony z kształtów.

A jednak z ułożenia ich razem wyszły trzy rzeczy, których żaden z tych
programów nie mógłby powiedzieć sam: że rzutowania zostawiają argument o
wariancji \emph{dokładnie} nienaruszony, że połączenie rezydualne zmienia to,
czym gradient \emph{jest}, a nie jak duży jest, i że warstwa gęsta mieści dwa
razy tyle, co uwaga.

\textbf{Tym właśnie jest architektura.} Nie nowym pomysłem, tylko kolejnością,
w jakiej stosuje się pomysły znane, dobraną tak, by założenia każdego z nich
były wciąż prawdziwe, gdy potrzebuje ich następny. Czytanie jej w ten sposób
było tym, po co ta książka powstała.
\end{fr}

\begin{summarybox}
\sumitem{1--2}{Blok to sześć operacji i ta książka zdefiniowała już wszystkie
  sześć. Architekturę tworzy kolejność i kształty.}
\sumitem{3--5}{\result{Parametry uwagi to cztery rzutowania; wynik to jeden iloczyn
  skalarny; a $n$ pozycji daje $n^{2}$ wyników}, co jest najbardziej brzemienną
  liczbą w tym projekcie.}
\sumitem{7--8}{Dzielnik wyniku to poprawka wariancji Programu~\ref{prog:P25}.
  W bloku zapytania i klucze są \emph{wyliczane}, więc pytaniem jest, czy ten
  argument przeżywa, a Program~\ref{prog:P25} przekazał je tutaj.}
\sumitem{9--10}{Przeżywa, dokładnie. \result{Wynik jest formą dwuliniową $x\T M y$ z
  $M = W_Q W_K\T$, jego wariancja to $\lVert M\rVert_F^{2}$, a to ma wartość
  oczekiwaną $d_k$} \dash{} więc rzutowania nic nie zmieniają i do stwierdzenia
  tego nie było potrzebne żadne losowanie.}
\sumitem{11--12}{\result{Ale wytrenowany model ma jedno losowanie wag, a nie średnią,
  więc jego własne wyniki są o kilka procent obok, zanim trening się zacznie}, a
  rozrzut maleje wraz z poszerzaniem głowicy, bo rząd $M$ jest co najwyżej
  $d_k$ (Program~\ref{prog:P08}).}
\sumitem{13--16}{\result{$\softmax$ czyni wyjście kombinacją wypukłą
  (Program~\ref{prog:P19}), więc warstwa czysto uwagowa nie może wyprodukować
  kierunku, którego jej wartości już nie rozpinają}.}
\sumitem{18--20}{\result{Uwaga wielogłowicowa jest darmowa: $\val{p32.heads}$ głowic
  szerokości $\val{p32.d.head}$ mieści dokładnie parametry jednej głowicy
  szerokości $\val{p32.d}$. Zmienia się rząd, a nie koszt}.}
\sumitem{21--22}{\result{Zwykły łańcuch $\val{p32.depth}$ warstw ma gradient co
  najwyżej $\val{p32.plain.bound}$, i to w jego \emph{najlepszym} przypadku.
  Liczba jest wciąż reprezentowalna \dash{} najmniejsza liczba normalna
  $\code{float32}$ to $\val{p01.fp32.minnorm}$ \dash{} więc zwykły łańcuch
  zabija bezużyteczność tego ograniczenia, a nie niedomiar. Niedomiar to
  przypadek realistyczny: przy warstwie nasyconej ten sam iloczyn spada poniżej
  $10^{-\val{f12.sat.exponent}}$, pod każdą podłogę tego formatu}.}
\sumitem{23--25}{\result{$y = x + f(x)$ czyni jakobian $\prod(1 + f'_k)$, którego
  rozwinięcie jest sumą po $2^{n}$ ścieżkach \dash{} po jednej na podzbiór
  warstw. \textbf{Pusty podzbiór wnosi dokładnie $1$.}} Gradient zwykłego
  łańcucha to jeden iloczyn; stosu rezydualnego \dash{} suma z $1$ w
  środku.}
\sumitem{26--27}{Normalizacja przywraca skalę, którą zakłada argument o
  wyniku, i dlatego te trzy kawałki są jednym mechanizmem.}
\sumitem{28}{\result{Obrót czyni wynik zależnym od odległości między dwiema pozycjami
  i od niczego więcej, a nic model nie kosztuje, bo jest ortogonalny}
  (Programy~\ref{prog:F08} i \ref{prog:P09}).}
\sumitem{29--32}{\result{Uwaga to $4d^{2}$, a warstwa gęsta to $8d^{2}$: \textbf{połowa,
  o której nikt nie mówi, mieści dwa razy tyle, co ta, o której mówią
  wszyscy}. Blok to $12d^{2}$}, a $\val{p32.layers}$ takich to
  $\val{p32.blocks.pct}$ procent tego, co Program~\ref{prog:F01} podaje
  jako $\val{p32.f01.params}$.}
\sumitem{33--35}{\result{Cache to jeden klucz i jedna wartość na pozycję na warstwę,
  czyli $\val{p32.kv.per.token.mib}$ MiB na token}, i odtwarza to, co
  Program~\ref{prog:P03} podaje jako $\val{p32.kv.gib}$ GiB. Podwój
  sekwencję: wyniki
  rosną czterokrotnie, cache dwukrotnie.}
\sumitem{36--39}{Każde wyprowadzenie tej książki o uwadze obowiązuje przy
  inicjalizacji, a pięć z nich opiera się na hipotezach, o które nie zapytano
  żadnego wytrenowanego modelu. Lista jest wydrukowana, a nie domyślna.}
\end{summarybox}

\canyou

\begin{testexercises}
\item Blok ma szerokość $d = 1024$. Ile parametrów jest w jego rzutowaniach
  uwagi, a ile w jego warstwie gęstej?
  \answerto{$4d^{2} = 4\,194\,304$ oraz $8d^{2} = 8\,388\,608$. Warstwa gęsta
  mieści dwa razy tyle, przy każdej szerokości \dash{} stosunek nie zależy od
  $d$.}
\item Model ma $16$ warstw szerokości $2048$. Oszacuj liczbę jego parametrów z
  samych bloków.
  \answerto{$12d^{2}L = 12 \times 2048^{2} \times 16$, czyli około
  $805$ milionów. Zanurzenie i głowica wyjściowa dochodzą na wierzch.}
\item Dwa zapytania i dwa klucze powstają przez rzutowania $W_Q$ i $W_K$.
  Które z $W_Q$, $W_K$ i $W_Q W_K\T$ da się odtworzyć z samych wyników?
  \answerto{Tylko iloczyn. Wyniki to $x\T W_Q W_K\T y$, więc każda para
  macierzy o tym samym iloczynie daje identyczne wyniki \dash{} i dlatego te
  dwie często przechowuje się i analizuje jako jedną macierz.}
\item Stos rezydualny ma $n$ warstw i każde $f'_k$ wypada dokładnie $-1$. Ile
  wynosi gradient i czy ścieżka identycznościowa go ratuje?
  \answerto{$\prod(1 + (-1)) = 0$. Ścieżka identycznościowa wnosi $1$, a
  pozostałe $2^{n} - 1$ ścieżek znosi to dokładnie. \enquote{Jeden składnik
  jest równy $1$} nie znaczy \enquote{suma jest co najmniej $1$}, i ten
  przypadek to pokazuje.}
\item Budżet obsługi pozwala na $32$ GiB cache'u przy kształtach z
  sekcji~\ref{sec:P32-count}. Ile rozmów po $\val{p32.seq}$ tokenów się
  mieści?
  \answerto{Osiem, bo każda to $\val{p32.kv.gib}$ GiB. Zmniejsz kontekst o
  połowę, a zmieści się szesnaście \dash{} cache jest liniowy w sekwencji, więc
  ta wymiana jest jeden do jednego.}
\item Która wielkość w bloku rośnie kwadratowo wraz z długością sekwencji, a
  która liniowo?
  \answerto{Liczba wyników rośnie kwadratowo; cache rośnie liniowo. Obie są
  własnościami jednej warstwy, i dlatego plan pojemności musi nazwać, której z
  nich zabrakło.}
\end{testexercises}


\begin{furtherproblems}
\item Pokaż, że $\Ex\lVert W_Q W_K\T\rVert_F^{2} = d_k$, gdy wyrazy obu
  macierzy mają wariancję $1/d$, licząc składniki.
  \answerto{$M$ ma $d^{2}$ wyrazów; każdy jest sumą $d_k$ iloczynów dwóch
  niezależnych wyrazów, więc każdy ma wariancję $d_k/d^{2}$; a
  $d^{2} \times d_k/d^{2} = d_k$. Szerokość strumienia się skraca, i dlatego
  dzielnik zależy od szerokości głowicy, a nie od szerokości modelu.}
\item Głowicę zwęża się, a liczbę głowic podnosi, by utrzymać stałą szerokość
  łączną. Co dzieje się z liczbą parametrów, a co z rzędem $M$ każdej głowicy?
  \answerto{Liczba parametrów się nie zmienia. Każde $M$ ma rząd co najwyżej
  $d_k$, który maleje, więc te same parametry wyrażają bardziej ograniczony
  zbiór porównań.}
\item Weź rozwinięcie $\prod_{k=1}^{n}(1 + f'_k)$ i zapisz sumę składników
  przechodzących dokładnie przez jedną warstwę.
  \answerto{$\sum_k f'_k$. Razem z $1$ ze ścieżki identycznościowej daje to
  przybliżenie pierwszego rzędu $1 + \sum_k f'_k$, i dlatego płytki stos
  rezydualny zachowuje się jak suma swoich warstw, a nie jak ich iloczyn.}
\item Wyników jest $n^{2}$, a cache jest liniowy w $n$. Przy jakiej długości
  sekwencji macierz wyników, po dwa bajty na liczbę i przy jednej głowicy,
  zrówna się z cache'em z sekcji~\ref{sec:P32-count}?
  \answerto{Z równania $2n^{2} = 2 \times \val{p32.layers} \times \val{p32.d}
  \times 2 \times n$ wychodzi $n = 2 \times \val{p32.layers} \times
  \val{p32.d}$, czyli $262\,144$ \dash{} daleko poza długości kontekstu
  używane tutaj, i dlatego przy obsługiwanych sekwencjach dominuje cache, a
  arytmetyce i tak dominują wyniki.}
\item Program~\ref{prog:F12} pokazuje, że zwykły łańcuch może też
  \emph{eksplodować}, a sekcja~\ref{sec:P32-residual} mówi, że ścieżka
  rezydualna tego nie naprawia. Dlaczego nie?
  \answerto{Ścieżka identycznościowa ogranicza gradient od dołu, a nie od góry.
  Pozostałe ścieżki są nietknięte, więc stos o dużych czynnikach wciąż je
  mnoży \dash{} po to jest normalizacja z sekcji~\ref{sec:P32-norm} i dlatego
  obie pojawiają się razem w każdym bloku.}
\end{furtherproblems}
